% \documentclass[aspectratio=169,notes]{beamer}
\documentclass[aspectratio=169]{beamer}
\usetheme[faculty=phil]{fibeamer}
\usepackage{polyglossia}
\setmainlanguage{english} %% main locale instead of `english`, you
%% can typeset the presentation in either Czech or Slovak,
%% respectively.
\setotherlanguages{russian} %% The additional keys allow
%%
%%   \begin{otherlanguage}{czech}   ... \end{otherlanguage}
%%   \begin{otherlanguage}{slovak}  ... \end{otherlanguage}
%%
%% These macros specify information about the presentation
\title[Artificial Intelligence Software (AIS)]{Artificial Intelligence Software, HW 1} %% that will be typeset on the
\subtitle{Reproucible environment setup:  
\\ Conda, UV, Docker \\ \  } %% title page.
\author{Oleg Bulichev}
%% These additional packages are used within the document:
\usepackage{ragged2e}  % `\justifying` text
\usepackage{booktabs}  % Tables
\usepackage{tabularx}
\usepackage{tikz}      % Diagrams
\usetikzlibrary{calc, shapes, backgrounds}
\usetikzlibrary{decorations.pathreplacing,calligraphy,calc,graphs}
\usepackage{amsmath, amssymb}
\usepackage{url}       % `\url`s
\usepackage{listings}  % Code listings
% \usepackage{subfigure}
\usepackage{floatrow}
\usepackage{subcaption}
\usepackage{mathtools}
\usepackage{todonotes}
\usepackage{fontspec}
\usepackage{multicol}
\usepackage{pdfpages}
\usepackage{wrapfig}
\usepackage{qrcode}
\usepackage{animate}
\usepackage{booktabs}
\usepackage{multirow}
% \usepackage{graphicx}
\usepackage{colortbl}

\graphicspath{{resources/}}
\frenchspacing

\setbeamertemplate{caption}[numbered]
\usetikzlibrary{graphs}

% \usepackage[backend=biber,style=ieee,autocite=footnote]{biblatex}
% \addbibresource{biblio.bib}
% \DefineBibliographyStrings{english}{%
%   bibliography = {References},}

\newcommand{\oleg}[2][] {\todo[color=red, #1] {OLEG:\\ #2}}
\newcommand{\fbckg}[1]{\usebackgroundtemplate{\includegraphics[width=\paperwidth]{#1}}}%frame background

\usepackage[framemethod=TikZ]{mdframed}
\newcommand{\dbox}[1]{
\begin{mdframed}[roundcorner=3pt, backgroundcolor=yellow, linewidth=0]
\vspace{1mm}
{#1}
\vspace{1mm}
\end{mdframed}
}

\begin{document}
\setlength{\abovedisplayskip}{0pt}
\setlength{\belowdisplayskip}{0pt}
\setlength{\abovedisplayshortskip}{0pt}
\setlength{\belowdisplayshortskip}{0pt}

\fbckg{fibeamer/figs/title_page.png}
\frame[c]{\setcounter{framenumber}{0}
    \usebeamerfont{title}%
    \usebeamercolor[fg]{title}%
    \begin{minipage}[b][6.5\baselineskip][b]{\textwidth}%
        \textcolor{black}{\raggedright\inserttitle}
    \end{minipage}
    % \vskip-1.5\baselineskip

    \usebeamerfont{subtitle}%
    \usebeamercolor[fg]{framesubtitle}%
    \begin{minipage}[b][3\baselineskip][b]{\textwidth}
        \raggedright%
        \insertsubtitle%
    \end{minipage}
    \vskip.25\baselineskip
}
%   \frame[c]{\maketitle}

\fbckg{fibeamer/figs/common.png}

\note{\scriptsize \begin{itemize}
        \item \
    \end{itemize}}

\begin{frame}[t]{Task 1}
\vspace{-0.3cm}
\textbf{Short description}:
Describe your experience with vibecoding (AI coding assistants). 

If you have no prior experience, describe your expectations or the reasons why you haven't tried it yet (instead of answering the Requirements section below).

\textbf{Artifacts}: Report in pdf format.

\textbf{Requirements}
\footnotesize
\begin{enumerate}
    \item What tasks did you try to solve using vibecoding? Which tools did you use?
    \item What did you like and dislike about the tools?
    \item What challenges or hallucinations did you encounter?
    \item How did you verify the correctness of the generated code?
    \item Which tools would you recommend to a friend and why?
    \item What would you like to learn about vibecoding in the next lecture?
\end{enumerate}
\end{frame}

\begin{frame}[t]{Task 2}
\vspace{-0.3cm}
\textbf{Short description}:
Build a Docker image with ROS~2 Humble based on \texttt{ubuntu:jammy} with GPU, GUI, and host file access support.

\textbf{Artifacts}: \texttt{Dockerfile}, \texttt{docker-compose.yml}, pdf file with brief explanation of your solution.


\textbf{Requirements}
\footnotesize
\begin{multicols}{2}
\begin{itemize}
  \item Base image: \texttt{ubuntu:jammy}
  \item ROS~2 Humble via official ROS repositories
  \item UID/GID mapping for host file access
  \item Start container using \texttt{docker compose}
  \item Pass \texttt{ROS\_DOMAIN\_ID} as environment variable
  \item Mounted volume for verification
  \item \texttt{network\_mode: host}
  \item NVIDIA Container Toolkit support (\texttt{nvidia-smi})
  \item GUI support (activate \texttt{rviz2})
  \item Container stays alive after \texttt{docker compose up} (can be solved via $command,\ tty\ or\ entrypoint$)
  \item Interactive access via \texttt{docker exec -it}
\end{itemize}
\end{multicols}

\end{frame}

\begin{frame}[t]{Task 3}
\framesubtitle{}
\vspace{-0.3cm}
\textbf{Short description}:
set up and compare a Data Science environment using 4 different tools: \textbf{Conda}, \textbf{Poetry}, \textbf{uv}, and \textbf{Pixi}.

\textbf{Artifacts}: Repository with configuration folders for each tool, \texttt{benchmark.py} script, and report in pdf format with comparison. Compare \textbf{Cold Install Time}, \textbf{Disk Usage}, and developer experience.

\textbf{Requirements}
\footnotesize
\begin{enumerate}
    \item \textbf{Benchmark Script}: Create \texttt{benchmark.py} that uses \texttt{numpy, pandas, matplotlib, plotly} to generate data and save a plot.
    \item \textbf{Conda}: Create env from \texttt{environment.yaml}, export full lock-file.
    \item \textbf{Poetry}: Init project, resolve dependencies, run script via \texttt{poetry run}.
    \item \textbf{uv}: Create venv, install deps. Focus on resolution speed.
    \item \textbf{Pixi}: Init project. Install Python deps + one system tool (e.g. \texttt{ffmpeg} or \texttt{graphviz}).
\end{enumerate}

\end{frame}

\fbckg{fibeamer/figs/last_page.png}
\frame[plain]{}

\end{document}